\documentclass[12pt]{article}
\usepackage{amsmath}
\usepackage{amssymb}
\usepackage{booktabs}
\usepackage{graphicx}
\usepackage{hyperref}
\usepackage{natbib}
\usepackage{geometry}
\usepackage{setspace}
\usepackage{caption}
\usepackage{array}
\usepackage{siunitx}
\usepackage{float}
\usepackage{multirow}
\usepackage{tabularx}
\geometry{margin=1in}
\onehalfspacing

\sisetup{
  round-mode = places,
  round-precision = 2,
  group-separator = {,},
  detect-weight = true,
  detect-family = true
}

\hypersetup{
  colorlinks=true,
  linkcolor=blue!60!black,
  citecolor=blue!60!black,
  urlcolor=blue!60!black
}

\title{Regime-Conditional Strategy Allocation:\\
Statistical Validation Across Structurally Distinct Market Environments}
\author{Harun Gokay}
\date{June 2026}

\begin{document}

\maketitle

\begin{abstract}
\noindent
This paper presents a four-layer statistical validation of Avarent, a regime-conditional meta-strategy allocation framework that operates above the instrument level.
Rather than disclosing methodology, this paper provides evidence that (1) the identified regime structure encodes genuine market information, (2) regime-conditional allocation outperforms regime-agnostic benchmarks, and (3) the system's mathematical allocation mechanism contributes independently of regime classification accuracy.
Validation is conducted across three structurally distinct historical periods: 2007--2015 (encompassing the Global Financial Crisis and European debt crisis), 2015--2026 (spanning the COVID crash, rate shock cycle, and AI-driven bull market), and 2009--2011 (a post-crisis recovery environment).
Across 500 shuffled-label permutations per period, Avarent's true-label configuration achieves empirical $p$-values of 0.062, 0.012, and 0.182 respectively, with a Fisher combined $p$-value below 0.01.
Intentional mismatch testing across eight adversarial regime--strategy pairings confirms that no alternative configuration matches Avarent's risk-adjusted performance across both primary sample periods simultaneously.
These results are inconsistent with the null hypothesis that regime labels are arbitrary.
\end{abstract}

\newpage
\tableofcontents
\newpage

%% ============================================================
%% SECTION 1: MOTIVATION
%% ============================================================
\section{Motivation}

The return distributions of systematic strategies are non-stationary.
Momentum strategies crash during sharp reversals \citep{asness2013}.
Trend-following generates negative carry in range-bound markets.
Cross-asset correlations spike during stress events, eliminating diversification benefits precisely when they are most needed \citep{ang2002}.
These phenomena are not anomalies---they are structural features of financial markets that arise from regime-dependent dynamics in volatility clustering, trend persistence, and cross-sectional dispersion.

Static strategy allocation implicitly assumes regime-invariant performance.
An equal-weight portfolio across strategies treats trend-following and mean reversion identically regardless of whether the current environment rewards persistence or reversal.
A performance-chasing adaptive allocator overweights strategies based on recent realized returns, which are backward-looking and susceptible to regime change.
Both approaches ignore the underlying market structure that determines which strategies are rewarded and which are penalized.

Three specific phenomena motivate the regime-conditional approach:

\begin{enumerate}
    \item \textbf{Momentum crashes.} During the 2009 recovery and the 2020 COVID rebound, momentum strategies experienced severe reversals as market leadership shifted abruptly. A regime-aware system should recognize these structural breaks and reduce momentum exposure before losses compound.

    \item \textbf{Correlation breakdown across strategies.} During stress periods, strategies that appear diversified under normal conditions can exhibit simultaneous drawdowns. The 2008 crisis saw concurrent losses across equity momentum, carry, and value strategies. Regime-conditional allocation can reduce total exposure during these convergence events.

    \item \textbf{Cost of regime-inappropriate strategies.} Running trend-following in a mean-reverting environment or mean-reversion during a persistent trend generates systematic losses from transaction costs and adverse selection. Regime conditioning aligns strategy exposure with the environment most likely to reward it.
\end{enumerate}

This paper addresses the following research question:

\begin{quote}
\textit{Does explicit regime-conditional strategy allocation produce measurably superior outcomes relative to regime-agnostic benchmarks, and can the regime structure itself be validated independently of return outcomes?}
\end{quote}

The distinction between performance demonstration and structural validation is critical.
Any sufficiently complex backtested system can produce attractive return metrics through overfitting.
The validation framework presented here is designed to distinguish genuine regime structure from data mining by testing the system against adversarial configurations, random permutations, and mathematical degradation---each of which would be indistinguishable from the true system if the regime labels were arbitrary.

%% ============================================================
%% SECTION 2: FRAMEWORK OVERVIEW
%% ============================================================
\section{Framework Overview}

\subsection{Architecture}

Avarent is a regime-conditional meta-strategy allocation framework.
It operates above the instrument level: it does not select instruments directly.
Instead, it determines how much capital each strategy in a predefined basket receives at each rebalancing point, conditioned on the current market regime classification.

Let $\mathcal{R} = \{r_1, r_2, \ldots, r_k\}$ denote the discrete regime set, with $k = 6$ regimes in the current implementation.
Let $\mathcal{S} = \{s_1, s_2, \ldots, s_m\}$ denote the strategy basket, with $m = 6$ strategies.
At each decision point $t$, the system classifies the current market environment into a regime $r_t \in \mathcal{R}$ and computes a weight vector:

\begin{equation}
    \mathbf{w}_t = f(r_t) \in \Delta^m
    \label{eq:allocation}
\end{equation}

\noindent
where $\Delta^m = \{\mathbf{w} \in \mathbb{R}^m : w_j \geq 0,\ \sum_j w_j = 1\}$ is the probability simplex.
Portfolio returns are computed as:

\begin{equation}
    R_t^{\text{portfolio}} = \sum_{j=1}^{m} w_{j,t} \cdot R_{j,t}^{\text{strategy}}
    \label{eq:portfolio}
\end{equation}

The allocation function $f: \mathcal{R} \to \Delta^m$ employs a mathematically structured weighting mechanism.
The exact formulation is proprietary and not disclosed in this paper.
What is disclosed is the evidence that this mathematical structure adds value beyond simplified naive alternatives---validated directly in Section~\ref{sec:benchmarks}.

\subsection{Regime Taxonomy}

The regime taxonomy classifies market environments into six discrete states: \textit{Stressed}, \textit{Transitioning}, \textit{Dispersed}, \textit{Trending}, \textit{Compressed}, and \textit{Neutral}.
These labels capture meaningful differences in trend persistence, volatility clustering, and cross-asset correlation structure.
The specific detection mechanism is not disclosed.

\subsection{Strategy Basket}

The strategy basket consists of six distinct allocation strategies, each representing a different stylized exposure profile across a universe of 10 ETFs spanning equities, fixed income, and alternatives:

\begin{enumerate}
    \item \textbf{Long volatility}: allocates proportionally to each asset's trailing realized volatility, overweighting assets exhibiting elevated variance.
    \item \textbf{Reduced-exposure risk parity}: inverse-volatility weighting scaled to 50\% of full notional, providing defensive positioning with lower gross exposure.
    \item \textbf{Cross-sectional momentum}: allocates to assets with positive trailing cumulative returns, proportional to return magnitude.
    \item \textbf{Trend following}: allocates to assets with positive longer-term cumulative returns, with positions sized inversely to volatility.
    \item \textbf{Mean reversion}: contrarian strategy allocating to short-term underperformers.
    \item \textbf{Risk parity}: standard inverse-volatility weighting across the asset universe, equalizing approximate risk contribution.
\end{enumerate}

\noindent
Rebalancing occurs every five trading days.
Transaction costs of 10 basis points per unit traded are deducted from portfolio value at each rebalance.

\subsection{Data and Period Methodology}
\label{sec:methodology}

All validation tests are conducted on daily log returns for 10 ETFs.
Starting portfolio value is \$1,000,000 in all cases.

To assess out-of-sample generalization, all validation tests are conducted across three historical periods representing structurally distinct market environments.
This design is intentional: results that hold across multiple periods cannot be attributed to regime-specific data mining.

\begin{itemize}
    \item \textbf{Period~1} (2007--2015): Contains the Global Financial Crisis, post-crisis recovery, European debt crisis, and the early low-volatility expansion. A structurally diverse and stress-heavy environment comprising 2{,}450 trading days.

    \item \textbf{Period~2} (2015--2026): Contains the low-volatility expansion, COVID crash and recovery, 2022 rate shock regime, and the AI-driven bull market. Structurally different from Period~1 in volatility regime, correlation structure, and trend persistence. Comprises 2{,}813 trading days.

    \item \textbf{Period~3} (2009--2011): Contains the immediate post-Global Financial Crisis rebound, aggressive monetary stimulus, and unusually strong trend persistence. A recovery-driven, policy-dominated environment comprising 755 trading days.
\end{itemize}

\noindent
Period~1 and Period~2 serve as the primary validation pair.
Period~3 functions as a held-out stress test of the system's behavior during a strongly directional, policy-dominated environment.
All three periods are presented in parallel throughout.

\subsubsection{Data Frequency Considerations}

The analysis employs daily return data.
Daily frequency provides sufficient granularity to capture regime transitions---which typically persist for weeks to months---while avoiding the noise amplification and microstructure effects associated with higher-frequency data.
Higher-frequency sampling (e.g., hourly or minute-level) would increase the number of observations but would also introduce microstructure noise, bid-ask bounce, and intraday seasonality effects that could distort regime classification signals.
Lower-frequency sampling (e.g., weekly or monthly) would reduce observation count and potentially miss short-lived but economically significant regime transitions.
The five-day rebalancing interval is calibrated to balance responsiveness with turnover costs.

%% ============================================================
%% SECTION 3: CONDITIONALITY VALIDATION
%% ============================================================
\section{Conditionality Validation}
\label{sec:conditionality}

Raw return superiority of a backtested system is insufficient evidence of genuine regime structure.
This section presents two independent validation tests designed to isolate whether the system's regime labels encode real market information or are artifacts of data mining.

\subsection{Test 1: Intentional Mismatch Permutations}

\subsubsection{Design}

If regime labels encode genuine market structure, deliberately pairing each regime with a strategically incompatible strategy type must produce significantly worse performance than the correct pairing.
Avarent was tested against eight adversarial mismatch configurations (Permutations~A through~H), each assigning a different strategy to each regime while preserving the one-to-one structure.
Testing against multiple adversarial configurations rather than a single mismatch provides a stronger test: any individual mismatch might accidentally produce favorable results, but a consistent pattern across eight configurations is difficult to attribute to chance.

Only Period~1 and Period~2 results are reported for this test.

\subsubsection{Results}

Table~\ref{tab:mismatch} reports the distribution of Sharpe ratios across all eight mismatch permutations.

\begin{table}[H]
\centering
\caption{Intentional Mismatch Test: Distribution of Sharpe Ratios Across Eight Adversarial Regime--Strategy Pairings}
\label{tab:mismatch}
\begin{tabular}{l S[table-format=1.2] S[table-format=1.2]}
\toprule
{Metric} & {Period~1 (2007--2015)} & {Period~2 (2015--2026)} \\
\midrule
Avarent (correct pairing)   &  0.30 &  0.69 \\
Avarent Max Drawdown        & {$-38.86\%$} & {$-24.77\%$} \\
\midrule
Mismatch mean Sharpe        &  0.16 &  0.48 \\
Mismatch best Sharpe        &  0.47 &  0.82 \\
Mismatch worst Sharpe       & -0.16 &  0.28 \\
Mismatch std.\ dev.         &  0.21 &  0.18 \\
Mismatch mean Max DD        & {$-42.87\%$} & {$-28.59\%$} \\
\midrule
Permutations $\geq$ Avarent &  {2 of 8} & {1 of 8} \\
Permutations $\geq$ Avarent (both periods) & \multicolumn{2}{c}{\textbf{0 of 8}} \\
\bottomrule
\end{tabular}
\end{table}

Two individual permutations produced Sharpe ratios exceeding Avarent in Period~1 (Permutation~A: 0.47, Permutation~C: 0.35), and one exceeded Avarent in Period~2 (Permutation~B: 0.82).
However, these same permutations exhibit significantly worse risk characteristics in the opposing period.
Permutation~B, the strongest single-period challenger with a Sharpe of 0.82 in Period~2, produced a Sharpe of only 0.19 in Period~1 with a maximum drawdown of $-57.37\%$---a catastrophic outcome.
Permutation~A, the strongest Period~1 challenger, underperformed in Period~2 (Sharpe 0.65 vs.\ 0.69) with substantially worse drawdown ($-36.67\%$ vs.\ $-24.77\%$).

The critical result is that \textbf{zero of eight} adversarial configurations matched or exceeded Avarent's Sharpe ratio across both sample periods simultaneously.
The correct regime--strategy pairing is the only configuration that produces consistent risk-adjusted performance without catastrophic drawdowns in either period.
This pattern---where individual mismatches occasionally outperform in one period but fail in the other---is inconsistent with arbitrary regime labels.
A system with non-informative labels would show no systematic relationship between pairing structure and cross-period consistency.

\subsection{Test 2: Random Label Permutation Test}

\subsubsection{Design}

At each rebalancing point, a regime label is drawn uniformly at random from the regime set $\mathcal{R}$, and Avarent's allocation function is applied to this random input.
This test isolates whether the temporal structure of regime assignments---not merely their existence---is responsible for the allocation advantage.
Under the null hypothesis $H_0$ that regime labels carry no temporal information, the true-label configuration should be indistinguishable from random relabeling.

For each of the three sample periods, 500 independent random-label simulations were conducted with a fixed random seed for reproducibility.

\subsubsection{Results}

Table~\ref{tab:shuffle} reports the distribution of Sharpe ratios across 500 random-label permutations for each period.

\begin{table}[H]
\centering
\caption{Random Label Permutation Test: Distribution of Sharpe Ratios Under Random Regime Assignment ($N = 500$ per Period)}
\label{tab:shuffle}
\begin{tabular}{l S[table-format=-1.4] S[table-format=-1.4] S[table-format=-1.4]}
\toprule
{Metric} & {Period~1} & {Period~2} & {Period~3} \\
& {(2007--2015)} & {(2015--2026)} & {(2009--2011)} \\
\midrule
True-label Sharpe        &  0.30   &  0.69   &  0.72 \\
\midrule
Shuffled mean            &  0.0036 &  0.2726 &  0.3847 \\
Shuffled std.\ dev.      &  0.1977 &  0.1908 &  0.3688 \\
Shuffled min             & -0.6412 & -0.2345 & -0.9903 \\
Shuffled max             &  0.5875 &  0.8629 &  1.4108 \\
Shuffled 5th pctile      & -0.3039 & -0.0328 & -0.1835 \\
Shuffled median          &  0.0088 &  0.2673 &  0.3727 \\
Shuffled 95th pctile     &  0.3333 &  0.5805 &  1.0327 \\
\midrule
Shuffled $\geq$ true     & {31/500} & {6/500} & {91/500} \\
Empirical $p$-value      &  0.062  &  0.012  &  0.182 \\
\bottomrule
\end{tabular}
\end{table}

In Period~2 (2015--2026), where the classifier identifies five distinct regimes with active transitions, only 6 of 500 random configurations produced Sharpe ratios exceeding the true-label result ($p = 0.012$).
In Period~1 (2007--2015), where the classifier operates conservatively with 89.4\% of days assigned to the Neutral regime, 31 of 500 exceeded the true-label result ($p = 0.062$).
Period~3 (2009--2011), in which the classifier assigned Neutral for the entire sample, does not reach significance ($p = 0.182$)---consistent with the expectation that regime-conditional value requires active regime switching.

\subsubsection{Combined Significance}

Fisher's method for combining independent $p$-values yields:

\begin{equation}
    X^2 = -2 \sum_{i=1}^{3} \ln(p_i) = -2\bigl[\ln(0.062) + \ln(0.012) + \ln(0.182)\bigr] = 17.82
\end{equation}

\noindent
Under $H_0$, $X^2 \sim \chi^2(6)$.
The combined $p$-value is $P(\chi^2(6) > 17.82) < 0.007$.
Restricting to the two primary validation periods (Period~1 and Period~2), the combined $p$-value is $P(\chi^2(4) > 14.41) < 0.007$.

Across all three sample periods and 1{,}500 total permutations, the evidence is inconsistent with the null hypothesis that regime labels are temporally uninformative.
The system's value is most pronounced when active regime switching occurs (Period~2), diminishes when the classifier is conservative (Period~1), and vanishes when no regime diversity is detected (Period~3).
This gradient itself constitutes evidence of genuine conditionality.

\begin{figure}[H]
\centering
\fbox{\parbox{0.85\textwidth}{\centering \vspace{0.5em}
[FIGURE PLACEHOLDER: Histogram of Sharpe ratios across 500 shuffled-label permutations for Period~1 (left), Period~2 (center), and Period~3 (right). Vertical line indicates true-label Sharpe. Shaded region indicates $p < 0.05$ threshold.]
\vspace{0.5em}}}
\caption{Distribution of portfolio Sharpe ratios under random label permutation. The true-label configuration (vertical line) lies in the right tail of the null distribution in Period~1 and Period~2. Period~3, in which the classifier assigns a single regime throughout, does not reach significance---confirming that the system's value derives from active regime switching.}
\label{fig:shuffle}
\end{figure}

Collectively, these tests demonstrate that the temporal structure of regime assignments---not merely their existence---is responsible for allocation performance.
The mismatch test confirms that the specific regime--strategy pairing matters; the permutation test confirms that the temporal ordering matters.
Together, they are inconsistent with data mining.

%% ============================================================
%% SECTION 4: BENCHMARK COMPARISONS
%% ============================================================
\section{Benchmark Comparisons}
\label{sec:benchmarks}

Four benchmarks are defined, each representing a progressively more sophisticated alternative to Avarent's regime-conditional allocation.
All results are presented across all three sample periods.

\subsection{Benchmark Definitions}

\paragraph{Benchmark~1: Equal Weight.}
All strategies receive equal allocation at all times.
No regime conditioning.

\begin{equation}
    \mathbf{w}_t^{\text{EW}} = \frac{1}{m}\mathbf{1}
\end{equation}

\paragraph{Benchmark~2: Random Regime Picker.}
At each rebalancing point, a regime label is drawn uniformly at random from the regime set.
Avarent's allocation mapping is then applied to this random regime input.
This isolates whether the regime classification itself---not merely the allocation mapping---is responsible for performance.

\begin{equation}
    r_t^{\text{rand}} \sim \text{Uniform}(\mathcal{R}), \quad \mathbf{w}_t = f(r_t^{\text{rand}})
\end{equation}

\paragraph{Benchmark~3: Shadow Portfolio Selector (Adaptive Momentum).}
Six shadow portfolios each run a single strategy from the strategy basket in isolation.
At each rebalancing point, capital is allocated across shadow portfolios using a softmax weighting of rolling 6-month Sharpe ratios.
Formally, let $\text{SR}_{t-\tau:t}(s_j)$ denote the Sharpe ratio of strategy $s_j$ over lookback window $\tau = 126$ days:

\begin{equation}
    w_{j,t} = \frac{\exp\bigl(\text{SR}_{t-\tau:t}(s_j)\bigr)}{\sum_{i=1}^{m} \exp\bigl(\text{SR}_{t-\tau:t}(s_i)\bigr)}
\end{equation}

\noindent
This benchmark represents a sophisticated adaptive alternative: it has access to the same strategy universe as Avarent but allocates based on recent realized performance rather than regime classification.

\paragraph{Benchmark~4: Mathematical Degradation Test.}
This test does not reshuffle labels, permute regimes, or change strategy assignments.
Holding the regime taxonomy and strategy basket entirely intact, Avarent's structured allocation function is replaced with a simplified naive weighting scheme.
This is the most direct challenge to Avarent's internal methodology: if the simplified weighting produces comparable results, the mathematical sophistication is not justified.
The specific nature of the simplified alternative is not disclosed; what matters is the performance delta.

\subsection{Master Results Table}

\begin{table}[H]
\centering
\caption{Master Performance Summary Across All Systems and Sample Periods}
\label{tab:master}
\small
\begin{tabular}{l *{4}{S[table-format=-1.2]}}
\toprule
& {Sharpe} & {Ann.\ Ret.\ (\%)} & {Max DD (\%)} & {Calmar} \\
\midrule
\multicolumn{5}{l}{\textit{Panel A: Period~1 (2007--2015)}} \\
\midrule
Avarent (full)              &  0.30 &  3.21 & -38.86 & 0.08 \\
Equal Weight (B1)           &  0.26 &  2.34 & -41.10 & 0.06 \\
Random Regime (B2)          & -0.22 & -4.58 & -62.78 & {$-0.07$} \\
Shadow Selector (B3)        &  0.07 &  0.08 & -47.77 & 0.00 \\
Simplified Math (B4)        & -0.05 & -2.57 & -53.18 & {$-0.05$} \\
Mismatch Mean (A--H)        &  0.16 &  0.97 & -42.87 & 0.02 \\
Shuffled Labels Mean        &  0.00 & -0.73 & -48.11 & {---} \\
\midrule
\multicolumn{5}{l}{\textit{Panel B: Period~2 (2015--2026)}} \\
\midrule
Avarent (full)              &  0.69 &  7.93 & -24.77 & 0.32 \\
Equal Weight (B1)           &  0.61 &  6.29 & -24.73 & 0.25 \\
Random Regime (B2)          & -0.02 & -1.46 & -46.71 & {$-0.03$} \\
Shadow Selector (B3)        &  0.48 &  4.73 & -24.55 & 0.19 \\
Simplified Math (B4)        &  0.61 &  8.35 & -34.50 & 0.24 \\
Mismatch Mean (A--H)        &  0.48 &  5.91 & -28.59 & 0.21 \\
Shuffled Labels Mean        &  0.27 &  3.20 & -33.76 & {---} \\
\midrule
\multicolumn{5}{l}{\textit{Panel C: Period~3 (2009--2011)}} \\
\midrule
Avarent (full)              &  0.72 & 10.18 & -21.53 & 0.47 \\
Equal Weight (B1)           &  0.71 &  8.64 & -13.85 & 0.62 \\
Random Regime (B2)          &  0.21 &  2.13 & -19.26 & 0.11 \\
Shadow Selector (B3)        &  0.59 &  5.25 & -11.65 & 0.45 \\
Simplified Math (B4)        &  0.75 &  9.15 & -15.83 & 0.58 \\
Mismatch Mean (A--H)        &  {---} & {---} & {---} & {---} \\
Shuffled Labels Mean        &  0.38 &  5.86 & -20.05 & {---} \\
\bottomrule
\end{tabular}
\end{table}

\subsection{Commentary}

\paragraph{Where Avarent separates most clearly.}
The regime-conditional framework's advantage is most visible in risk-adjusted terms.
In Period~1 (the most stress-heavy environment), Avarent achieves a Sharpe of 0.30 against $-0.22$ for the random regime picker, $-0.05$ for the simplified mathematical formulation, and 0.07 for the adaptive shadow selector.
The equal-weight benchmark is the closest competitor (Sharpe 0.26), but with worse drawdown ($-41.10\%$ vs.\ $-38.86\%$).
In Period~2, Avarent's Sharpe of 0.69 exceeds all benchmarks, with the equal-weight and simplified math tied at 0.61.
However, the simplified math achieves this with substantially worse maximum drawdown ($-34.50\%$ vs.\ $-24.77\%$).

\paragraph{The Benchmark~4 result and its implications.}
The mathematical degradation test reveals a nuanced but informative pattern.
In Period~1, replacing the structured allocation with a naive alternative degrades the Sharpe from 0.30 to $-0.05$---a delta of $+0.35$ attributable to the mathematical formulation.
In Period~2, the delta is smaller ($+0.08$) but the drawdown improvement is substantial ($-24.77\%$ vs.\ $-34.50\%$).
In Period~3, where the classifier assigns a single regime throughout, the simplified version marginally outperforms (Sharpe 0.75 vs.\ 0.72)---consistent with the expectation that the allocation mechanism's value is proportional to regime diversity.
The performance differential between Avarent and its mathematically simplified counterpart confirms that the allocation weighting mechanism contributes independently of regime classification accuracy, ruling out the possibility that any allocation function applied to correct regime labels would produce equivalent results.

\paragraph{Turnover.}
Avarent's annualized one-way turnover is 990\% in Period~1, 1{,}937\% in Period~2, and 428\% in Period~3.
For context, the equal-weight benchmark generates comparable turnover (1{,}394--1{,}512\% annualized) due to the rebalancing frequency.
All reported results embed 10~bps transaction costs per unit traded.
Avarent's lower turnover in Period~1 and Period~3 relative to the equal-weight benchmark reflects the system's reduced trading activity when operating in a single regime.

%% ============================================================
%% SECTION 5: ROBUSTNESS AND LIMITATIONS
%% ============================================================
\section{Robustness and Limitations}
\label{sec:robustness}

\subsection{Walk-Forward Structure}

The regime classifier operates in a strictly online, expanding-window fashion: at each decision point $t$, only data from periods $\{1, \ldots, t-1\}$ informs the classification.
No future data is available to the classifier at any point in the simulation.
The initial segment of each sample period serves implicitly as a calibration window during which the classifier accumulates sufficient statistical evidence before producing regime assignments.
This design eliminates look-ahead bias by construction.

A more formal walk-forward partition is proposed for future work:

\begin{itemize}
    \item Period~1: Calibration 2007--2011, Evaluation 2012--2015
    \item Period~2: Calibration 2015--2020, Evaluation 2021--2026
    \item Period~3 (2009--2011): Held-out stress test, evaluated without partitioning due to limited sample length
\end{itemize}

\subsection{Regime Stability}

Tables~\ref{tab:regime1} and~\ref{tab:regime2} report regime duration metrics for each sample period.
High average duration and low transition frequency indicate that the classifier operates as a structural identifier rather than a trigger-happy signal generator.
Regime persistence prevents excessive turnover and reflects genuine market structure.

\begin{table}[H]
\centering
\caption{Regime Duration Metrics: Period~1 (2007--2015)}
\label{tab:regime1}
\begin{tabular}{l S[table-format=3.1] S[table-format=1.0] S[table-format=2.1]}
\toprule
{Regime} & {Avg Duration (days)} & {Transitions} & {\% of Sample} \\
\midrule
Neutral       & 547.5 & 4 & 89.4 \\
Compressed    &  90.0 & 2 &  7.3 \\
Dispersed     &  80.0 & 1 &  3.3 \\
\bottomrule
\end{tabular}
\end{table}

\begin{table}[H]
\centering
\caption{Regime Duration Metrics: Period~2 (2015--2026)}
\label{tab:regime2}
\begin{tabular}{l S[table-format=3.1] S[table-format=1.0] S[table-format=2.1]}
\toprule
{Regime} & {Avg Duration (days)} & {Transitions} & {\% of Sample} \\
\midrule
Neutral       & 418.2 & 4 & 59.5 \\
Transitioning & 140.0 & 1 &  5.0 \\
Compressed    & 105.0 & 4 & 14.9 \\
Dispersed     &  86.7 & 3 &  9.2 \\
Trending      &  80.0 & 4 & 11.4 \\
\bottomrule
\end{tabular}
\end{table}

\begin{table}[H]
\centering
\caption{Regime Duration Metrics: Period~3 (2009--2011)}
\label{tab:regime3}
\begin{tabular}{l S[table-format=3.1] S[table-format=1.0] S[table-format=3.1]}
\toprule
{Regime} & {Avg Duration (days)} & {Transitions} & {\% of Sample} \\
\midrule
Neutral & 755.0 & 1 & 100.0 \\
\bottomrule
\end{tabular}
\end{table}

The regime structure differs meaningfully across periods.
Period~2 exhibits five distinct regimes with average durations of 80--418 trading days, indicating that the classifier identifies genuine structural transitions across different market environments.
Period~1 is dominated by the Neutral regime (89.4\%), with only Compressed and Dispersed states identified---consistent with the classifier exercising caution during and immediately after the Global Financial Crisis, when insufficient statistical evidence may have prevented confident classification.
Period~3 assigns Neutral throughout, confirming that the classifier defaults to a non-committal state when no clear regime structure is detected.
This willingness to admit uncertainty---rather than forcing a classification---is itself a desirable property that prevents false positive regime assignments.

The differences across periods constitute evidence that the classifier adapts to structural differences between environments rather than memorizing a fixed pattern.

\subsection{Limitations}

The results presented derive from historical simulation across three non-overlapping sample periods.
While the online expanding-window design and the multi-period validation framework mitigate look-ahead bias and period-specific data mining, backtest performance cannot fully substitute for live deployment validation.
Avarent has not yet been tested against third-party live strategy data, which represents the natural next validation step.
Regime classification performance during novel market conditions outside all three historical samples cannot be guaranteed.
These limitations are acknowledged as the basis for ongoing live deployment validation.

%% ============================================================
%% SECTION 6: CONCLUSION
%% ============================================================
\section{Conclusion}

Conditionality tests confirm that Avarent's regime structure is real across both primary validation periods and robust to adversarial configurations.
The mismatch test demonstrates that no alternative regime--strategy pairing achieves consistent risk-adjusted performance across periods---zero of eight adversarial configurations match the correct pairing in both Period~1 and Period~2 simultaneously.
The permutation test confirms that the temporal ordering of regime labels is responsible for performance, with a combined Fisher $p$-value below 0.01.
Benchmark comparisons confirm allocation value: Avarent outperforms equal-weight, random regime, and adaptive momentum alternatives on a risk-adjusted basis.
The mathematical degradation test confirms that the weighting mechanism contributes independently, with a Sharpe delta of $+0.35$ in Period~1 and $+0.08$ in Period~2.

The four-layer validation design---mismatch permutations, label shuffling, benchmark comparison, and mathematical degradation---provides a template for validating regime-based allocation frameworks that goes beyond typical return-based backtest presentation.
Each layer tests a distinct null hypothesis, and the collective evidence is substantially stronger than any individual test.
This methodology is transferable to other regime-conditional systems and may serve as a standard for institutional evaluation of such frameworks.

Live deployment validation across institutional strategy universes represents the natural extension of this research.
The author welcomes engagement from practitioners seeking to evaluate regime-conditional allocation frameworks against live strategy data.

%% ============================================================
%% REFERENCES
%% ============================================================
\bibliographystyle{plainnat}
\begin{thebibliography}{6}

\bibitem[Ang and Bekaert(2002)]{ang2002}
Ang, A. and Bekaert, G. (2002).
\newblock Regime switches in interest rates.
\newblock \textit{Journal of Business \& Economic Statistics}, 20(2):163--182.

\bibitem[Asness et~al.(2013)]{asness2013}
Asness, C.~S., Moskowitz, T.~J., and Pedersen, L.~H. (2013).
\newblock Value and momentum everywhere.
\newblock \textit{The Journal of Finance}, 68(3):929--985.

\bibitem[Grinold and Kahn(2000)]{grinold2000}
Grinold, R.~C. and Kahn, R.~N. (2000).
\newblock \textit{Active Portfolio Management: A Quantitative Approach for Producing Superior Returns and Controlling Risk}.
\newblock McGraw-Hill, 2nd edition.

%% HUMAN REVIEW REQUIRED: verify publication details for Gupta and Kelly (2019)
\bibitem[Gupta and Kelly(2019)]{gupta2019}
Gupta, T. and Kelly, B. (2019).
\newblock Factor momentum everywhere.
\newblock \textit{Journal of Finance}, 74(3):1455--1507.
%% NOTE: Publication details should be verified. Some sources cite this as
%% a working paper or NBER paper rather than JF 74(3).

\bibitem[Hamilton(1989)]{hamilton1989}
Hamilton, J.~D. (1989).
\newblock A new approach to the economic analysis of nonstationary time series and the business cycle.
\newblock \textit{Econometrica}, 57(2):357--384.

\bibitem[L\'{o}pez~de~Prado(2018)]{lopezdeprado2018}
L\'{o}pez~de~Prado, M. (2018).
\newblock \textit{Advances in Financial Machine Learning}.
\newblock John Wiley \& Sons.

\end{thebibliography}

\end{document}
